\section{Đánh giá thiết kế}

\subsection{Mục tiêu đánh giá}
% Mục tiêu kiểm thử:
% - Kiểm tra tính dễ dùng và hiệu quả của luồng đặt lịch mới
% - Đo lường mức độ rõ ràng của tính năng theo dõi tiến độ và dặn dò dị ứng

\subsection{Phương pháp đánh giá}
% Phương pháp kiểm thử:
% - Kiểm thử tính khả dụng (Usability Testing) kết hợp giao thức Think-aloud
% - Bảng câu hỏi đo lường mức độ hài lòng (Thang đo tính khả dụng hệ thống - SUS)

\subsection{Người tham gia đánh giá}
% Số lượng và tiêu chí tuyển chọn người tham gia kiểm thử (ví dụ: 5 chủ nuôi thú cưng)

\subsection{Tác vụ kiểm thử}
% Danh sách các tác vụ kiểm thử (gắn liền với các kịch bản trong Proposal):
% - Tác vụ 1: Đặt lịch tắm & cắt tỉa cho thú cưng có tiền sử dị ứng
% - Tác vụ 2: Kiểm tra và theo dõi tiến độ chăm sóc trên ứng dụng
% - Tác vụ 3: Xem lại lịch sử dịch vụ và ghi chú của lần chăm sóc trước

\subsection{Chỉ số đo lường}
% Các chỉ số đo lường:
% - Tỷ lệ hoàn thành tác vụ (Task Completion Rate - %)
% - Thời gian thực hiện tác vụ (Time on Task - giây)
% - Tỷ lệ sai sót (Error Rate) / Số lỗi phát sinh
% - Điểm số SUS (SUS Score)

\subsection{Quy trình đánh giá}
% Mô tả quy trình tổ chức buổi kiểm thử

\subsection{Kết quả đánh giá}
% Trình bày kết quả đo lường thực tế qua bảng biểu và biểu đồ

\subsection{Thảo luận và phân tích}
% Thảo luận ý nghĩa kết quả:
% - Những điểm cải thiện vượt trội so với quy trình cũ
% - Những vấn đề còn tồn đọng người dùng gặp phải trong quá trình test

